\section{Reward Modeling}
\label{sec:reward-modeling}

A reward model (RM) replaces a hand-written objective with a learned scoring function.  For a prompt $x$ and completion $y$, it estimates a scalar or a set
of local scores that can be used for best-of-$N$ selection, search, or policy optimization.  Unlike an environment reward, this signal is inferred from human preferences or correctness annotations and is therefore only a proxy for the desired behavior.

This section uses $x$ for a prompt, $y$ for a completion, $y_c$ for a chosen completion, and $y_r$ for a rejected completion.  It is adapted from Lambert's reward-modeling chapter.\footnote{
\url{https://rlhfbook.com/c/05-reward-models}}

\subsection{Bradley--Terry Preference Reward Model}

\subsubsection{Preference likelihood}

The Bradley--Terry (BT) model assigns each item $i$ a positive latent strength
$p_i$.  Its pairwise preference probability is
\begin{equation}
    P(i \succ j)
    =
    \frac{p_i}{p_i+p_j}.
\end{equation}
Writing $p_i=\exp(r_i)$ gives the logit parameterization
\begin{equation}
    P(i \succ j)
    =
    \frac{\exp(r_i)}{\exp(r_i)+\exp(r_j)}
    =
    \sigma(r_i-r_j),
    \qquad
    \sigma(z)=\frac{1}{1+\exp(-z)}.
\end{equation}
Consequently, only score differences are identifiable: replacing every score
$r_i$ by $r_i+C$ leaves all preference probabilities unchanged.

For an RM $r_\theta(x,y)$ and a preference triple
$(x,y_c,y_r)\sim\mathcal{D}$,
\begin{equation}
    P_\theta(y_c\succ y_r\mid x)
    =
    \sigma\!\left(
        r_\theta(x,y_c)-r_\theta(x,y_r)
    \right).
    \label{eq:bt-probability}
\end{equation}
Maximum likelihood yields the final loss function:
\begin{align}
    \mathcal{L}_{\mathrm{BT}}(\theta)
    &=
    \mathbb{E}_{(x,y_c,y_r)\sim\mathcal{D}}
    - \log \sigma\!\left(
        r_\theta(x,y_c)-r_\theta(x,y_r)
    \right).
    % &=
    % \mathbb{E}_{(x,y_c,y_r)\sim\mathcal{D}}
    % \operatorname{softplus}\!\left(
    %     r_\theta(x,y_r)-r_\theta(x,y_c)
    % \right).
    \label{eq:bt-loss}
\end{align}
The log is applied before averaging; maximizing
$\mathbb{E}[P_\theta(y_c\succ y_r\mid x)]$ is not the same likelihood
objective.

\subsubsection{Architecture}

The standard architecture reuses a causal-transformer backbone and replaces
its vocabulary projection with a scalar head.  For prompt-completion tokens
$z=(x,y)$, let $h_t\in\mathbb{R}^d$ be the final-layer hidden state.  A common
sequence score is
\begin{equation}
    r_\theta(x,y)=w^\top h_{t_{\mathrm{last}}}+b,
\end{equation}
where $t_{\mathrm{last}}$ is the final non-padding token, normally an EOS
token.  Chosen and rejected sequences are scored independently by the same
network.  The loss in \eqref{eq:bt-loss} couples only their scalar outputs.

\begin{codeblock}[language=Python,caption={Minimal Bradley--Terry reward model.}]
import torch
import torch.nn as nn
import torch.nn.functional as F

class BradleyTerryRM(nn.Module):
    def __init__(self, backbone):
        super().__init__()
        # Transformer producing one d-dimensional hidden state per token.
        self.backbone = backbone
        # Shared linear head: final token representation -> scalar reward.
        self.reward_head = nn.Linear(backbone.config.hidden_size, 1)

    def forward(self, input_ids, attention_mask):
        """
        Score each complete prompt-response sequence.

        input_ids:      token IDs, shape [batch_size, sequence_length]
        attention_mask: 1 for real tokens and 0 for padding, same shape
        returns:        one unnormalized reward logit per sequence, shape [B]
        """
        out = self.backbone(
            input_ids=input_ids,
            attention_mask=attention_mask,
            output_hidden_states=True,
            return_dict=True,
        )

        # Last transformer layer: one contextual representation per token.
        hidden = out.hidden_states[-1]              # [B, T, d]

        # Locate the final non-padding token. Unlike mask.sum() - 1, this is
        # correct for both left-padded and right-padded batches.
        positions = torch.arange(
            hidden.size(1), device=hidden.device
        ).expand_as(attention_mask)
        last = positions.masked_fill(attention_mask == 0, -1).max(dim=1).values
        batch = torch.arange(hidden.size(0), device=hidden.device)
        final_hidden = hidden[batch, last]           # [B, d], usually EOS

        # The score is not a probability; only score differences enter BT loss.
        return self.reward_head(final_hidden).squeeze(-1)  # [B]

def preference_loss(model, chosen, rejected):
    """Negative log-likelihood that chosen is preferred to rejected."""
    # Both dictionaries contain the same prompts but different completions.
    # The same reward model (shared parameters) scores both sequences.
    reward_chosen = model(**chosen)
    reward_rejected = model(**rejected)

    # -log sigmoid(r_chosen - r_rejected), averaged over the batch.
    # Minimization pushes the chosen score above the rejected score.
    return -F.logsigmoid(reward_chosen - reward_rejected).mean()
\end{codeblock}

Correct construction of the two inputs is critical: each contains the same
prompt and exactly one completion, and the attention mask must select the
actual final token under either left or right padding.  Reward models are
commonly trained for only one epoch because pairwise datasets are correlated
and overfitting can rapidly degrade held-out ranking accuracy.

\subsection{Outcome Reward Models}

An outcome reward model (ORM) predicts whether a completion eventually reaches
a correct final answer.  It is trained from independently labeled examples
$(x,y,r)$, where $r\in\{0,1\}$ denotes final-answer correctness; paired
chosen--rejected data is unnecessary.

One common ORM emits a binary logit $w_{\theta,t}$ at every completion token.
The sequence outcome label is broadcast across completion positions while
prompt and padding positions are masked:
\begin{equation}
    \mathcal{L}_{\mathrm{ORM}}(\theta)
    =
    -\mathbb{E}_{(x,y,r)\sim\mathcal{D}}
    \left[
        \frac{1}{\sum_t m_t}
        \sum_{t=1}^{T}m_t
        \left(
            r\log p_{\theta,t}
            +(1-r)\log(1-p_{\theta,t})
        \right)
    \right],
    \label{eq:orm-token-loss}
\end{equation}
where $m_t=1$ only for completion tokens and
$p_{\theta,t}=\sigma(w_{\theta,t})$.  These targets contain no direct
information about whether an intermediate step is valid: every token receives
the label determined by the final outcome.

A sequence-pooled alternative first averages logits and then applies the
sigmoid,
\begin{equation}
    \bar p_\theta(x,y)
    =
    \sigma\!\left(
        \frac{1}{\sum_t m_t}\sum_t m_t w_{\theta,t}
    \right),
\end{equation}
followed by sequence-level binary cross entropy.  In general,
$\sigma(\operatorname{mean}(w_t))\neq
\operatorname{mean}(\sigma(w_t))$, so these aggregation orders are not
interchangeable.

\begin{codeblock}[language=Python,caption={Per-token outcome reward model.}]
import torch.nn as nn
import torch.nn.functional as F

class OutcomeRM(nn.Module):
    def __init__(self, backbone):
        super().__init__()
        self.backbone = backbone
        self.outcome_head = nn.Linear(backbone.config.hidden_size, 1)

    def forward(self, input_ids, attention_mask, labels=None):
        out = self.backbone(
            input_ids=input_ids,
            attention_mask=attention_mask,
            output_hidden_states=True,
            return_dict=True,
        )
        hidden = out.hidden_states[-1]
        logits = self.outcome_head(hidden).squeeze(-1)  # [B, T]

        if labels is None:
            return {"logits": logits}

        # labels: -100 for prompt/padding; 0 or 1 for completion tokens
        mask = labels.ne(-100)
        loss = F.binary_cross_entropy_with_logits(
            logits[mask], labels[mask].float()
        )
        return {"loss": loss, "logits": logits}
\end{codeblock}

At inference, the per-token probabilities can be reduced to a candidate score
using the mean, the last few positions, the minimum, or a sum of log
probabilities.  The reduction must match validation-time behavior; products
strongly penalize long completions and minimum aggregation is highly sensitive
to one anomalous token.

\subsection{Process Reward Models}

A process reward model (PRM) evaluates intermediate reasoning steps rather than
only the final answer.  Let
$s=(s_1,\ldots,s_K)$ be an annotated reasoning trace and let
$q_i\in\{0,1\}$ indicate whether step $s_i$ is valid conditioned on the prompt
and previous steps.  A binary PRM minimizes
\begin{equation}
    \mathcal{L}_{\mathrm{PRM}}(\theta)
    =
    -\mathbb{E}_{(x,s)\sim\mathcal{D}}
    \sum_{i=1}^{K}
    \left[
        q_i\log r_\theta(s_i\mid x,s_{<i})
        +(1-q_i)\log\!\left(1-r_\theta(s_i\mid x,s_{<i})\right)
    \right].
    \label{eq:prm-loss}
\end{equation}
In practice, a separator token marks each step boundary.  All non-boundary
positions receive the ignore index $-100$.  A common three-class formulation
uses class indices $0,1,2$ for incorrect, neutral, and correct.

\begin{codeblock}[language=Python,caption={Three-class process reward model.}]
import torch.nn as nn
import torch.nn.functional as F

class ProcessRM(nn.Module):
    def __init__(self, backbone, num_classes=3):
        super().__init__()
        self.backbone = backbone
        self.step_head = nn.Linear(
            backbone.config.hidden_size, num_classes
        )

    def forward(self, input_ids, attention_mask, labels=None):
        out = self.backbone(
            input_ids=input_ids,
            attention_mask=attention_mask,
            output_hidden_states=True,
            return_dict=True,
        )
        logits = self.step_head(out.hidden_states[-1])  # [B, T, 3]

        if labels is None:
            return {"logits": logits}

        # labels: 0=incorrect, 1=neutral, 2=correct at boundaries;
        #         -100 at every other token.
        loss = F.cross_entropy(
            logits.reshape(-1, logits.size(-1)),
            labels.reshape(-1),
            ignore_index=-100,
        )
        return {"loss": loss, "logits": logits}
\end{codeblock}

For search or reranking, PRM outputs are read only at step boundaries.  Step
scores can be averaged, summed with position-dependent weights, or reduced by
the minimum to reject a trace as soon as one step is judged invalid.

\subsection{Reward Models versus Value Functions}

The output shape alone does not define the model.  The supervision target and
its semantics are decisive:

\begin{center}
\small
\begin{tabularx}{\textwidth}{@{}l X X l@{}}
\toprule
Model & Prediction & Supervision & Output location \\
\midrule
BT RM
& Relative quality $r_\theta(x,y)$
& Pairwise preferences
& EOS/last token \\
ORM
& Probability of eventual correctness
& Fixed offline outcome label
& Every completion token \\
PRM
& Validity of the current reasoning step
& Step-level annotations
& Step boundaries \\
Value function
& Expected remaining return $V^\pi(s_t)$
& Return targets from policy rollouts
& Every policy token \\
\bottomrule
\end{tabularx}
\end{center}

A value function predicts
\begin{equation}
    V^\pi(s_t)
    =
    \mathbb{E}_{\pi}\!\left[
        \sum_{k=t}^{T}\gamma^{k-t}r_k
        \,\middle|\,s_t
    \right].
\end{equation}
Its targets generally depend on the current policy's rollouts and change as
that policy changes.  By contrast, an ORM is normally trained offline against
fixed correctness labels broadcast to token prefixes.  Even if both use the
same backbone and scalar per-token head, the ORM estimates outcome
correctness, whereas the value head estimates remaining return and serves as a
baseline for advantages such as
$A_t=\widehat{R}_t-V_\phi(s_t)$.

Training a sequence-level BT model on pairs formed by one correct and one
incorrect response does \emph{not} turn it into an ORM.  If the objective is
\eqref{eq:bt-loss} and the output is one scalar per sequence, it remains a
pairwise preference RM; correctness merely determines how the pairs were
constructed.

\subsection{Inference Interfaces}

\begin{description}
    \item[BT RM.]
    Input $(x,y)$; output one scalar $r_\theta(x,y)$.  Use it as a terminal
    RL reward or rank $N$ sampled completions by their scalar scores.

    \item[ORM.]
    Input $(x,y)$; output
    $p_t\approx P(\text{final answer correct}\mid x,y_{\leq t})$ at completion
    positions.  Aggregate token scores to verify or rerank finished
    candidates.

    \item[PRM.]
    Input $(x,s_1,\ldots,s_K)$; output correctness scores at step boundaries.
    Aggregate across steps for reranking, or use scores online to prune
    low-quality search branches.

    \item[Value function.]
    Input the current prefix state $s_t=(x,y_{<t})$; output $V_\phi(s_t)$.
    During RL training it provides a per-token baseline rather than an
    interpretable probability that the whole response is correct.
\end{description}
